\documentclass[12pt]{article}

% === Packages ===
\usepackage[margin=1in]{geometry}
\usepackage{amsmath, amssymb}
\usepackage{setspace}
\usepackage{titlesec}
\usepackage{enumitem}
\usepackage{helvet}
\renewcommand{\familydefault}{\sfdefault} % modern sans-serif look

% === Formatting tweaks ===
\setstretch{1.15}
\setlist[enumerate]{itemsep=0.4em, topsep=0.2em}
\titleformat{\section}{\Large\bfseries\sffamily}{\thesection.}{1em}{}
\titleformat{\subsection}{\large\bfseries\sffamily}{\thesubsection}{1em}{}

% === Title ===
\title{\vspace{-2em}\textbf{Applied Mathematics Oral Exam Questions and Model Answers}\\
\large Solid Mechanics}
\author{Prepared for Graduate Oral Examination}
\date{}
\begin{document}
\maketitle
%===========================================================
% Sample oral exam questions + sample answers
% Solid mechanics for IBFE (based on IBFE_Tutorials (14))
%===========================================================

\begin{enumerate}

%-----------------------------------------------------------
\item \textbf{Configurations and deformation map.}
State what the reference and current configurations are, and define the deformation map $\chi$.

\textbf{Sample answer:}
The reference (material) configuration is $\mathcal{B}_0$ with coordinates $X\in\mathcal{B}_0$.
The current (spatial) configuration is $\mathcal{B}_t$ with coordinates $x\in\mathcal{B}_t$.
The deformation (motion) is the map
\[
x=\chi(X,t):\mathcal{B}_0\to\mathcal{B}_t.
\]

%-----------------------------------------------------------
\item \textbf{Deformation gradient.}
Define the deformation gradient $F_\chi$ and give its geometric meaning.

\textbf{Sample answer:}
\[
F_\chi(X,t)=\nabla_X \chi(X,t).
\]
It is the local linear map sending reference line elements to current line elements:
\[
dx = F_\chi\, dX.
\]

%-----------------------------------------------------------
\item \textbf{Jacobian and volume change.}
Define $J$ and explain how it relates $dV$ and $dv$.

\textbf{Sample answer:}
\[
J(X,t)=\det(F_\chi(X,t)).
\]
It is the local volume ratio, so volume elements transform as
\[
dv = J\, dV.
\]
Physically admissible deformations typically require $J>0$.

%-----------------------------------------------------------
\item \textbf{Quick computation: isotropic stretch.}
If $\chi(X)=(\lambda X_1,\lambda X_2,\lambda X_3)$, compute $F_\chi$ and $J$.

\textbf{Sample answer:}
\[
F_\chi=\lambda I,\qquad J=\det(\lambda I)=\lambda^3.
\]

%-----------------------------------------------------------
\item \textbf{Quick computation: simple shear.}
If $\chi(X)=(X_1+\gamma X_2,\;X_2,\;X_3)$, compute $F_\chi$ and $J$.

\textbf{Sample answer:}
\[
F_\chi=
\begin{pmatrix}
1&\gamma&0\\
0&1&0\\
0&0&1
\end{pmatrix},
\qquad
J=\det(F_\chi)=1.
\]
So simple shear is volume-preserving.

%-----------------------------------------------------------
\item \textbf{Incompressibility condition.}
What does incompressibility mean in this framework?

\textbf{Sample answer:}
Local incompressibility means volume preservation:
\[
J(X,t)=1\quad \text{for all } X.
\]
Equivalently, $I_3=\det(C)=J^2=1$.

%-----------------------------------------------------------
\item \textbf{Right/left Cauchy--Green tensors.}
Define $C$ and $B$ and state what they represent.

\textbf{Sample answer:}
\[
C = F_\chi^T F_\chi \quad \text{(right Cauchy--Green, reference-frame)},
\qquad
B = F_\chi F_\chi^T \quad \text{(left Cauchy--Green, spatial-frame)}.
\]
They encode stretch information with rotations factored out.

%-----------------------------------------------------------
\item \textbf{Green--Lagrange strain.}
Define the Green--Lagrange strain tensor $E$.

\textbf{Sample answer:}
\[
E=\frac12(C-I)=\frac12(F_\chi^T F_\chi - I).
\]
It is a finite-strain measure defined in the reference configuration.

%-----------------------------------------------------------
\item \textbf{Small strain vs finite strain.}
State the small-strain tensor and when it is valid.

\textbf{Sample answer:}
With displacement $d(X,t)=\chi(X,t)-X$,
\[
\varepsilon=\frac12\bigl(\nabla_X d + (\nabla_X d)^T\bigr),
\]
which is valid when displacement gradients are small so higher-order terms can be neglected.
For larger deformations, use $E=\tfrac12(C-I)$ instead.

%-----------------------------------------------------------
\item \textbf{Strain invariants.}
State the standard invariants $I_1,I_2,I_3$ of $C$.

\textbf{Sample answer:}
\[
I_1=\mathrm{tr}(C),\qquad
I_2=\frac12\left[(\mathrm{tr}C)^2-\mathrm{tr}(C^2)\right],\qquad
I_3=\det(C)=J^2.
\]
$I_3$ encodes volume change through $J$.

%-----------------------------------------------------------
\item \textbf{Cauchy stress and traction.}
Define Cauchy stress $\sigma$ and traction $t$.

\textbf{Sample answer:}
The Cauchy stress $\sigma(x,t)$ is the spatial stress tensor.
The traction on a current surface with unit normal $n$ is
\[
t(x,t;n) = \sigma(x,t)\, n,
\]
and the differential force on area element $da$ is $df=\sigma n\, da$.

%-----------------------------------------------------------
\item \textbf{First Piola--Kirchhoff stress.}
Define the first Piola--Kirchhoff stress $P$ and interpret $P N$.

\textbf{Sample answer:}
$P(X,t)$ is a reference-frame stress measure such that the force on a reference surface patch
with unit normal $N$ and area $dA$ is
\[
df = P(X,t)\,N\, dA.
\]
So $P N$ is the reference traction (force per reference area), expressed as a spatial vector.
In general $P$ is not symmetric.

%-----------------------------------------------------------
\item \textbf{Nanson's formula.}
State Nanson's formula and explain why it matters.

\textbf{Sample answer:}
Nanson's formula relates oriented area vectors:
\[
n\, da = J\,F_\chi^{-T}\,N\, dA.
\]
It allows us to equate the same physical force written in the current and reference descriptions,
and it is the key step in deriving the Piola transform between $P$ and $\sigma$.

%-----------------------------------------------------------
\item \textbf{Piola transform.}
Derive (at a high level) the relationship between $P$ and $\sigma$.

\textbf{Sample answer:}
Equate the same force on corresponding surface patches:
\[
\sigma n\, da = P N\, dA.
\]
Use Nanson's formula $n\, da = JF_\chi^{-T}N\, dA$ to get
\[
\sigma(JF_\chi^{-T}N\, dA)=PN\, dA.
\]
Since this holds for all $N$, we obtain
\[
P = J\,\sigma\,F_\chi^{-T}.
\]
Equivalently,
\[
\sigma = \frac{1}{J}P F_\chi^{T}.
\]

%-----------------------------------------------------------
\item \textbf{Second Piola--Kirchhoff stress.}
Define $S$ and relate it to $P$.

\textbf{Sample answer:}
\[
S = F_\chi^{-1}P,\qquad \text{equivalently } P = F_\chi S.
\]
$S$ is defined in the reference frame and is often symmetric for hyperelastic materials.

%-----------------------------------------------------------
\item \textbf{Why can $P$ be non-symmetric even if $\sigma$ is symmetric?}
\textbf{Sample answer:}
Even if $\sigma=\sigma^T$, the Piola transform $P=J\sigma F_\chi^{-T}$ multiplies $\sigma$
by $F_\chi^{-T}$, which generally destroys symmetry. $P$ is a two-point (mixed-index) tensor.

%-----------------------------------------------------------
\item \textbf{Momentum balance: reference form.}
State the strong form of momentum balance on $\mathcal{B}_0$.

\textbf{Sample answer:}
A standard reference form is
\[
\rho_0(X)\,\ddot{\chi}(X,t) = \nabla_X\cdot P(X,t) + B(X,t),
\]
where $\rho_0$ is reference density and $B$ is body force per reference volume.

%-----------------------------------------------------------
.


%-----------------------------------------------------------
\item \textbf{Why use invariants in isotropic materials?}
\textbf{Sample answer:}
For isotropic materials, the energy must be independent of rigid rotations.
Writing $W$ in terms of invariants of $C$ (or $B$), like $I_1,I_2,I_3$, guarantees rotational invariance.

%-----------------------------------------------------------
\item \textbf{Compute invariants for a diagonal stretch.}
If $F_\chi=\mathrm{diag}(\lambda_1,\lambda_2,\lambda_3)$, compute $J$, $I_1$, and $I_3$.

\textbf{Sample answer:}
\[
J=\lambda_1\lambda_2\lambda_3,\qquad
C=\mathrm{diag}(\lambda_1^2,\lambda_2^2,\lambda_3^2),
\]
so
\[
I_1=\lambda_1^2+\lambda_2^2+\lambda_3^2,\qquad
I_3=\det(C)=J^2=(\lambda_1\lambda_2\lambda_3)^2.
\]

%-----------------------------------------------------------
\item \textbf{Volumetric vs distortional response (concept check).}
What is meant by separating volumetric and distortional parts of the energy?

\textbf{Sample answer:}
One often writes $W=W_{\text{iso}}(\text{shape})+W_{\text{vol}}(J)$.
$W_{\text{vol}}(J)$ penalizes volume change (bulk response), while $W_{\text{iso}}$ controls shape change (shear-like response).
In the incompressible limit, one enforces $J=1$ (often via a pressure-like Lagrange multiplier).



\end{enumerate}


\end{document}